\part{Cross-readings}
\label{part:cross}

The ledger splits cleanly into campaigns.
Several of the open rows are the same formula read on different objects.
This part records those coincidences.
A local derivation is repeated, in shorter form, under the hypothesis it came from.

\section{One bond number}
\label{sec:bond}

\subsection{The identities}

Let \(\psi_i,\psi_j\in\C\) and write
\begin{equation}
\label{eq:Bq}
B_{ij}=\Re(\psi_i^*\psi_j),\qquad
q_{ij}=|\psi_i|^2+|\psi_j|^2,\qquad
s=\psi_i+\psi_j,\qquad
d=\psi_i-\psi_j.
\end{equation}
The sum-map convention in \texttt{split0.fiber\_point} is the same \(s\) and \(d\),
up to which daughter is called \(i\).
The swap \(d\mapsto -d\) is the endpoint gauge already quotiented by STORE.

\begin{lemma}
\label{lem:bond-expand}
\begin{align*}
|s|^2&=q_{ij}+2B_{ij},&
|d|^2&=q_{ij}-2B_{ij},\\
|s|^2-|d|^2&=4B_{ij},&
|s|^2+|d|^2&=2q_{ij}.
\end{align*}
If \(q_{ij}>0\),
\begin{equation}
\label{eq:Pminus}
P_-=\frac{|d|^2}{|s|^2+|d|^2}
=\frac12-\frac{B_{ij}}{q_{ij}},\qquad
P_+=\frac12+\frac{B_{ij}}{q_{ij}}.
\end{equation}
Furthermore \(|B_{ij}|\le q_{ij}/2\), with equality if and only if
\(\psi_i\) and \(\psi_j\) are nonnegative real multiples of one complex phase.
\end{lemma}

\begin{proof}
Expand \(|\psi_i\pm\psi_j|^2\).
The cross term is \(2\Re(\psi_i^*\psi_j)\).
Cauchy--Schwarz gives the bound on \(B_{ij}\).
The unitary modes \(\alpha=s/\sqrt{2}\) and \(\beta=d/\sqrt{2}\)
produce the same ratio \(P_-\), because the factors of \(2\) cancel.
\end{proof}

\begin{corollary}
\label{cor:Bzero}
For \(q_{ij}>0\),
\[
B_{ij}=0
\iff |s|=|d|
\iff P_+=P_-=\tfrac12.
\]
The sum-map norm change is \(\Delta\|\psi\|^2=2B_{ij}\)
(\texttt{accounting.event\_ledger}).
Hence \(B_{ij}=0\) if and only if that contraction is norm-neutral.
The energy change is
\(\Delta E_\psi=2B_{ij}-2\Sigma_{\mathrm{cross}}\).
On \(K_N\) every third vertex is a common neighbour, so
\(\Sigma_{\mathrm{cross}}=0\) and norm neutrality coincides with energy neutrality.
On a cut leg the exclusive-neighbour sum survives, and the two neutralities diverge.
\end{corollary}

If \(q_{ij}=0\) both amplitudes vanish.
There is then no pair, and \(P_-\) is undefined.
The limit of \(h_2(P_-)\) along \(q\to 0\) depends on the direction of approach,
so an empty pair is not assigned the value \(1\) by continuity.

\subsection{Binary entropy, as an optional reading}

\begin{remark}
\label{rem:h2}
INFO-0 recorded counts with binary logarithms and did not introduce
\(-\sum p\log p\).
The map \(B/q\mapsto P_-\) does not need a further identification.
You could still attach to an occupied pair the Shannon entropy of its
\(\pm\) weight,
\begin{equation}
\label{eq:se}
s_e=h_2(P_-),\qquad
h_2(p)=-p\log_2 p-(1-p)\log_2(1-p).
\end{equation}
Then \(0\le s_e\le 1\), with \(s_e=1\) if and only if \(B_{ij}=0\),
and \(s_e=0\) at the Cauchy--Schwarz ends.
This \(s_e\) is the uncertainty of one projective \(\pm\) measurement.
It is not the von Neumann entropy of a pure global state.
It does not see the phase of \(d\) relative to \(s\).
At \(B_{ij}=0\) that phase is a pure current:
\(|\Im(\psi_i^*\psi_j)|=|\psi_i||\psi_j|\).
The continuous fibre dimension \(\dim_\R d=2\) counts magnitude and phase.
The number \(s_e\) counts neither as a dimension.
The two are not proportional, which is the overlap between the
post-STORE entropy row and the BH-Q dimension count
(Section~\ref{sec:kernel}).
\end{remark}

\begin{lemma}
\label{lem:taylor}
With \(x=B/q\) and \(|x|\le 1/2\),
\[
h_2\Bigl(\tfrac12-x\Bigr)
=1-\frac{2}{\ln 2}\,x^2-\frac{4}{3\ln 2}\,x^4+O(x^6).
\]
There is no odd term.
\end{lemma}

\begin{proof}
Set \(f(x)=(\tfrac12-x)\ln(\tfrac12-x)+(\tfrac12+x)\ln(\tfrac12+x)\),
so \(h_2=-f/\ln 2\).
Then \(f(0)=-\ln 2\), \(f'(0)=0\),
\(f''(x)=4/(1-4x^2)\), \(f''(0)=4\), \(f'''(0)=0\), \(f''''(0)=32\).
The degree-\(4\) Taylor polynomial is \(-\ln 2+2x^2+\tfrac43 x^4\).
\end{proof}

\subsection{Three ensembles, one pair}

\begin{theorem}
\label{thm:haar}
Let \(\psi\) be Haar-uniform on the unit sphere in \(\C^N\), \(N\ge 2\),
and let \(\{i,j\}\) be any fixed pair.
Then \(P_-\) is uniform on \([0,1]\), and
\[
\mathbb{E}[h_2(P_-)]=\frac{1}{2\ln 2}.
\]
The expectation is exact at finite \(N\), not a large-\(N\) limit.
Weighting pairs by \(q_{ij}\) does not change it:
\(h_2(P_-)\) depends only on the fraction, and for the Dirichlet law
the fraction is independent of the pair norm.
\end{theorem}

\begin{proof}
The squared moduli of a Haar vector in \(\C^N\) are Dirichlet with every
parameter equal to \(1\).
They have the law of \((e_1,\ldots,e_N)/\sum_k e_k\) with \(e_k\) i.i.d.\
exponential.
For any two coordinates,
\(e_i/(e_i+e_j)\) is uniform on \([0,1]\).
The rotation
\(\alpha=(\psi_i+\psi_j)/\sqrt{2}\),
\(\beta=(\psi_i-\psi_j)/\sqrt{2}\)
is a global unitary.
Haar measure is invariant, so
\(P_-=|\beta|^2/(|\alpha|^2+|\beta|^2)\) is uniform on \([0,1]\).
The integral \(\int_0^1 -p\ln p\,dp=1/4\), and the partner term in \(p\mapsto 1-p\)
is equal, so
\[
\int_0^1 h_2(p)\,dp=\frac{1}{\ln 2}\cdot\frac12=\frac{1}{2\ln 2}.
\]
\end{proof}

\begin{proposition}
\label{prop:arcsine}
If the two moduli are equal and the relative phase is uniform,
or if the pair is a real Gaussian,
then \(P_-\) is Beta\(\bigl(\tfrac12,\tfrac12\bigr)\) and
\[
\mathbb{E}[h_2(P_-)]=2-\frac{1}{\ln 2}.
\]
\end{proposition}

\begin{proof}
Equal moduli and uniform phase give
\(P_-=(1-\cos\phi)/2\).
The density of this law is \(1/(\pi\sqrt{p(1-p)})\).
A real Gaussian pair has the same norm fraction, because the argument of
\((g_1,g_2)\) is uniform.
One of the two symmetric terms in the average is
\[
\frac{1}{\pi\ln 2}\int_0^1\frac{-p\ln p}{\sqrt{p(1-p)}}\,dp
=\frac{1}{\pi\ln 2}\bigl(-\partial_a B(a,b)\bigr)_{a=3/2,\,b=1/2}.
\]
Here \(B(3/2,1/2)=\pi/2\) and \(\psi(3/2)-\psi(2)=1-2\ln 2\),
so \(-\partial_a B=(\pi/2)(2\ln 2-1)\).
One term equals \((2\ln 2-1)/(2\ln 2)\).
Both terms give \(2-1/\ln 2\).
\end{proof}

\begin{proposition}
\label{prop:ground}
The normalised all-ones vector on \(K_N\), which is the ground state of
\(H=-A\), has \(B_{ij}/q_{ij}=1/2\) on every edge, hence \(P_-=0\) and \(h_2=0\).
\end{proposition}

A flat-magnitude random-phase vector is Proposition~\ref{prop:arcsine},
not Theorem~\ref{thm:haar}.
The two means differ:
\(1/(2\ln 2)\approx 0.721\) against \(2-1/\ln 2\approx 0.557\).

\subsection{Scale separation}

\begin{proposition}
\label{prop:lambda}
Fix \(v>0\) and let \(\psi_i=\sigma g\) with \(g\) complex normal,
\(\mathbb{E}|g|^2=1\), and \(\psi_j=v\) real.
Set \(\lambda=\sigma/v\). Then
\[
x=\frac{B}{q}=\frac{\lambda\,\Re g}{1+\lambda^2|g|^2}.
\]
As \(\lambda\to 0\) or \(\lambda\to\infty\), \(x\to 0\) in probability,
so \(h_2(1/2-x)\to 1\) in probability.
For small \(\lambda\), \(\mathbb{E}[(\Re g)^2]=1/2\) and
\[
1-\mathbb{E}[h_2]=\frac{\lambda^2}{\ln 2}+O(\lambda^4).
\]
\end{proposition}

You could read this as follows.
A dynamical selection of quadrature is one way to reach \(h_2=1\).
A mismatch of the two amplitudes is another, and it does not use the
phase. The case \(\psi_j=0\), \(\psi_i\neq 0\) is exactly \(B=0\).
That endpoint is the \(\psi=0\) limit, which the ledger already separates
from the JOINT vacua.

\subsection{Where the same number sits}

I1a assigns \(\ln 2\) nats, hence one bit, to a saturated leg, and the
matching theorem gives \(A=4\ln 2\cdot k\,\ell_P^2\).
In bits, \(S=k\) and
\[
\frac{S}{A}=\frac{1}{4\ell_P^2\ln 2}.
\]
You could take the channel set to be the \(k\) cut legs and set
\(S_Q=\sum_e h_2(P_{-,e})\).
Then \(S_Q=k\) if and only if \(\bar h=1\),
and the coefficient above is reproduced with no free area.
A Haar law on those same legs would give \(S_Q=k/(2\ln 2)\) instead,
short by the constant \(2\ln 2\).
The interior edges of \(K_N\) are a different set:
there are \(\binom{N}{2}\) of them, and Theorem~\ref{thm:haar}
puts about \(0.721\) bit on each.
Edges of \(K_N\) incident to a leg-bearing vertex number about \(kN\).
The cardinality of the set is part of the definition of the sum.
The algebra does not choose the set.

BR2-QUADRATURE says that structural response tracks \(\cos\Delta\theta\)
through \(B\) and a staggered flux tracks \(\sin\Delta\theta\) through the current.
It does not set \(B_{ij}=0\).
The word ``quadrature'' in that verdict and the condition \(B=0\)
share a trigonometric function and not a theorem.

\section{Kernel dimension is a different count}
\label{sec:kernel}

\begin{lemma}
\label{lem:fiber}
The sum map \((p,q)\mapsto s=p+q\) has kernel the complex line \(d=p-q\),
real dimension \(2\), for every \(s\) including \(s=0\).
After the \(U(1)\) quotient the physical dimension drops to \(1\) only for
the all-zero merged state.
The swap \(d\mapsto -d\) is discrete and does not halve the dimension.
\end{lemma}

A tree of \(n-1\) interior merges therefore carries naive continuous
dimension \(2(n-1)\), minus all-zero exceptions.
That is the count \(D_Q=2N_d+D_c\) in the ledger.
Lemma~\ref{lem:fiber} and Remark~\ref{rem:h2} do not convert it into \(h_2\).

\begin{proposition}
\label{prop:vol}
Suppose each merge contributes an independent complex coordinate,
and the exterior imprint of those coordinates factors through a complex
boundary space of dimension at most \(b=|\partial R|\).
Then
\[
D_Q\ge 2(N_Q-b).
\]
For a space-filling collapse, \(N_Q=n_R-1\), so the lower bound is
volume-leading.
If the exterior evolution is the collapsed graph alone, \(Q\) does not
enter \(H\) (the store is frozen between events), the Jacobian on the
\(d\)-coordinates vanishes, and \(D_Q=2(n_R-1)\) up to all-zero exceptions.
\end{proposition}

An area law for this kernel would need either \(N_Q\) itself of order \(b\),
or an imprint whose rank cancels the bulk and leaves a boundary layer.
Neither is implied by \(H=-A\) or by the sum map.
You could have a volume law for \(D_Q\) and, on a different set of edges,
an area law for \(\sum h_2\).
They count different things: interior relative modes, and leg channels.
Promoting the factor \(2\) in \(D_Q=2N_d\) to two bits per merge
is a third count, distinct from I1a's one bit per saturated leg.

\section{One length in the chain weave}
\label{sec:ell}

The chain construction places stitch endpoints at density \(\rho=4\lambda\)
on the sheets and sets
\begin{equation}
\label{eq:ell}
\ell_W=\rho^{-1/2}=\frac{1}{2\sqrt{\lambda}}
\end{equation}
in sheet cells.
A transverse step costs an in-plane walk of that typical length, plus one
stitch edge.

\begin{proposition}
\label{prop:metric}
In coordinates \((x,y,z)\) with \(z\) the sheet index, the Riemannian
approximation consistent with that cost is
\begin{equation}
\label{eq:ds}
ds^2=dx^2+dy^2+\ell_W^2\,dz^2.
\end{equation}
The linear axis ratio is \(\ell_W\).
In-plane shifts and a transpose act on \((x,y)\).
They do not rotate \(z\) into the sheet.
The quadratic symbol of \(J_2\) is already isotropic in the sheet
(\(\mathrm{Hess}=4I\)), so those moves have no quadratic in-plane
anisotropy left to erase.
\end{proposition}

Volume growth in the window \(\ell_W\ll r\ll S\ell_W\) is
\(V\sim r^3/\ell_W\) for a one-dimensional meeting graph.
An ellipsoid has the same volume exponent as a ball.
The exponent \(3\) and the ratio \(\ell_W\) are independent readings of
\eqref{eq:ds}.

\begin{proposition}
\label{prop:over}
\(\ell_W=1\) forces \(\lambda=1/4\).
The stub fraction \(1-\exp(-2\lambda)\) is then \(1-e^{-1/2}\approx 0.393\),
which the weave preregistration labels overconnected.
Every woven ladder value \(\lambda\le 0.08\) has \(\ell_W\ge 1.77\).
\end{proposition}

The same \(\ell_W\) is the crossover of \(d_{\mathrm{eff}}(r)\) and,
if stitches are later removed, the transverse factor in any scale
read off \eqref{eq:ds}.
A scalar \(a\propto 1/\rho_S\) built from this metric tracks \(a_\perp\),
not an isotropic scale factor.
The in-plane factor stays \(1\) while the sheet lattice is held fixed.
The ledger average \(M_{ab}=\langle n_a n_b\rangle\) needs normals in an
ambient three-space.
The chain construction does not have those normals.
An isotropic plane process would be a different ensemble, and
Proposition~\ref{prop:metric} would not apply to it.

On a finite ring the third direction is compact, circumference \(S\ell_W\).
For \(S=\infty\) the window is the infrared and the axis ratio remains.
Spectral dimension of the homogenized walk is \(3\) only inside the time
window corresponding to \eqref{eq:ell}, and returns to \(2\) after the ring
is explored.
That return is the compactification, not a second mechanism.

\section{One quartic ratio}
\label{sec:quartic}

At \(J=1\),
\begin{align}
\label{eq:disp}
\varepsilon_{J_2}(q)
&=-8+2|q|^2-\tfrac16(q_x^4+q_y^4)+O(q^6),\\
\varepsilon_{J_3}(q)
&=-12+2|q|^2-\tfrac16\sum_i q_i^4+O(q^6).
\end{align}
Both are sums of functions of one momentum component, so every
coefficient of \(q_i^2 q_j^2\) vanishes.
The Hessian at the band bottom is \(4I\).

\begin{theorem}
\label{thm:ratio}
Write the quartic part as
\(\beta\sum_i q_i^4+\gamma\sum_{i<j}q_i^2 q_j^2\).
A quartic function of \(|q|^2\), and the degree-\(4\) term in
\(\sqrt{m^2c^4+c^2p^2}\), both require \(\gamma/\beta=2\).
For \eqref{eq:disp}, \(\beta=-1/6\) and \(\gamma=0\).
No choice of \(m\), \(c\), or \(J\) removes the mismatch.
The first rotationally anisotropic correction and the first departure from
the relativistic series sit at the same order.
\end{theorem}

The long-wave ray direction can still be radial.
On \(J_2\),
\(v_i=4q_i-(2/3)q_i^3+\cdots\), so \(v\) is parallel to \(q\) on the axes
and the diagonals, and the deflection elsewhere is \(O(q^2)\).
The Manhattan cone is a different object: for
\(\varepsilon=-4\sum_{i=1}^d\cos q_i\),
\[
\frac{v_E^{\max}}{v_M}=\frac{1}{\sqrt{d}},
\]
equal to \(1/\sqrt{2}\) on \(J_2\) and \(1/\sqrt{3}\) on \(J_3\).
A threshold on the peak and a threshold in the forerunner read these two
speeds. The ratio does not require a continuum fit.

A chain weave has Hessian \(\mathrm{diag}(4,4,t_\perp)\) at quadratic order.
Rescaling the sheet-index coordinate matches this to \(4I_3\) and does not
match the quartic, nor the flat band.
The \(J_3\) zero eigenvalue is structural at every crystal momentum.
A stitch with \(\|P_- H P_+\|\neq 0\) lifts the corresponding sheet zero.
Nullity against stitch density is the spectral distinction;
a quadratic packet, after the rescaling, does not see it.

Dimension universality in the earned sense is universality of the
quadratic form: SCALE-0 has \(d_H\to 2\) with \(J_2\) and the square control
indistinguishable.
Fringes at a few cells see \(O(q^2)\) corrections and the plaquette.
Two graphs with the same Hessian and different \(\gamma/\beta\)
agree on \(d_H\) and need not agree on interference.
That is why a spectral feature vector can fail to classify F
while still classifying the volume exponent.

\section{Band edge, dimension, and a gap}
\label{sec:dos}

Near a band edge whose Hessian is a multiple of the identity,
\begin{equation}
\label{eq:dos}
\rho(E)\propto
\begin{cases}
(E-E_0)^{-1/2}& d=1,\\
\mathrm{const}& d=2,\\
(E-E_0)^{d/2-1}& d\ge 2.
\end{cases}
\end{equation}
The \(J_2\) dispersive edge is the \(d=2\) case.
With \(m^*=1/4\) and \(J=1\),
\(\rho\to m^*/(2\pi)=1/(8\pi)\) per quotient cell.
The \(J_3\) edge is the \(d=3\) case, vanishing as \(\sqrt{E-E_0}\).

Fermi's golden-rule width \(\Gamma=2\pi|V|^2\rho(E)\) therefore stays finite
as \(E\downarrow E_0\) on \(J_2\) and can be parametrically small on \(J_3\).
Dimension \(3\) is the lowest dimension in \eqref{eq:dos} with a vanishing
edge density, and it is also the weave value
\(d_{\mathrm{obs}}=d_{\mathrm{sheet}}+d_M\) for a \(2\)-dimensional sheet
and a \(1\)-dimensional meeting complex.
The density of states does not forbid \(d\ge 4\).
The weave count does, and only as a minimality reading:
P0' imports Maxwell dimension \(2\) when it sets \(\langle z\rangle=4\),
and the graph alone does not force that integer.

\begin{proposition}
\label{prop:clique}
On infinite \(J_2\) the essential spectrum of \(H=-A\) is \([-8,8]\) together
with the flat band at \(0\).
A finite edge perturbation does not move the essential spectrum.
If a vertex subset induces a clique \(K_n\) with \(n\ge 10\),
the normalised indicator of that subset has Rayleigh quotient \(-(n-1)\)
for \(H\), so
\[
E_0\le -(n-1)\le -9,\qquad m=-8-E_0\ge n-9.
\]
\end{proposition}

For \(n\le 9\) this trial vector does not clear the edge.
A triangle scores \(-2\), inside the band.
Modes inside the essential spectrum hybridize with \eqref{eq:dos}.
You could regard a motif with no eigenvalue below the edge as a resonance
of the earned unitary, and a motif with such an eigenvalue as a spectral
bound state.
The flat band on pure \(J_2\) is a protected kernel, extensive in the
sheet-antisymmetric sector, and it does not propagate in the quotient.

The event account \(R=A_{\mathrm{coef}}+|d|^2/2+\Re(\bar d\, W)\)
telescopes along a sequence of merges to the change in known energy.
Between events \(E_\psi\) is conserved.
Q-DYN-0b treats \(R\) as event-local, not as a term stored in \(Q\).
The spectral gap of Proposition~\ref{prop:clique} depends only on the
final graph and the final eigenmode.
Two histories that reach the same \((G,\psi)\) share that gap and need not
share \(\sum R\).

At fixed \(\psi\), an edge-set change shifts \(E_\psi\) only through
\(-2\sum(\Delta A_e) B_e\).
There is no extra term in \(H=-A\) coupling to \(d_{\mathrm{eff}}\).
Deleting stitches on which \(B_e=0\) changes ball growth and does not change
\(E_\psi\).
A mass read from the eigenvalue and a dimension read from the ball are
independent at that locus.

Isolated \(K_n\) has spectrum \(\{-(n-1)\}\) once and \(\{+1\}\) with
multiplicity \(n-1\).
Only the ground level can sit below \(-8\), and only for \(n\ge 10\).
The level \(+1\) lies in \([-8,8]\).

\section{Codimension along a unitary orbit}
\label{sec:codim}

Between events, \(i\partial_t\psi=H(G)\psi\) with \(G\) and \(Q\) fixed,
and \(\partial H/\partial Q=0\).
VACEXC0 gives the same carrier evolution on every JOINT background.
Arrival times read from \(G\), or from the linear packet, are constant on
each inter-event interval.
Any scale factor built from those times has \(\dot a=0\) there.
A frozen texture in the flat band does not move \(H\).
Orientations that are functions of \(G\) are likewise constant.
A deterministic drift of \(M_{ab}\) would be an update rule for \(G\),
which the present equations do not contain.
REWIRE0 leaves degree-preserving re-pairings tied, so there is also no
bond drift that would rotate a frame.

Unweighted edges have combinatorial length \(1\).
A non-integer dilation is not an automorphism of a fixed simple graph.

\begin{lemma}
\label{lem:mu}
Across a sum-map split let \(\mu\) be the difference between the split-side
sum \(p+q\) and the merged-side amplitude at the same vertex.
Then \(\mu(t_0)=0\) for every fibre point, and
\begin{equation}
\label{eq:mudot}
\dot\mu(t_0)=i\Bigl(s+\sum_{v\in C}\psi_v\Bigr),
\end{equation}
which does not depend on \(d=\psi_i-\psi_j\).
The \(d\)-dependent piece of \(\ddot\mu(t_0)\) is
\(-\tfrac12 d\bigl(|X_i|-|X_j|\bigr)\).
\end{lemma}

\begin{proof}
From \(i\partial_t\psi=-A\psi\) one has \(\partial_t\psi=iA\psi\).
On the split graph,
\[
\partial_t(p+q)=i\Bigl(s+\sum_{X_i}\psi+\sum_{X_j}\psi+2\sum_C\psi\Bigr).
\]
On the merged graph the exclusive and common neighbours each appear once.
Subtracting cancels the exclusive sums and leaves \eqref{eq:mudot}.
Differentiating the exclusive sums brings down \(|X_i|p\) and \(|X_j|q\),
and \(p-q=d\) produces the second-derivative coefficient.
\end{proof}

If \(s+\sum_C\psi\neq 0\), no choice of \(d\) restores \(C^1\) continuity of
the summed amplitude.
A uniform nonzero state fails \eqref{eq:mudot} on every cover.
A STORE replay meets every order because the recovered state is the
predecessor, not because \(d\) was tuned.
An asymmetric cover that has already passed \eqref{eq:mudot} fixes at most
one \(d\) at the next order. Higher orders then overdetermine it.

Equation~\eqref{eq:mudot} is two real constraints.
A real condition \(B_{ij}(t)=B_*\) is one.
Along the unitary orbit, \(B_{ij}(t)\) is almost periodic, with frequencies
equal to eigenvalue differences of \(H\).
For a two-mode parent
\(B(t)=B_0+C\cos(\omega t+\varphi)\),
crossings of a level inside the range exist and sit at an arccosine comb.
A one-dimensional orbit misses a codimension-\(2\) submanifold for generic
initial data, and meets a codimension-\(1\) level set in isolated roots.
The ratios inside \(B_*\) are unfixed (the conditional conservation identity
leaves two ratios free).
The family of roots, or of Rice rates
\[
\lambda(B_*)=\frac{1}{2\pi}\sqrt{\frac{m_2}{m_0}}
\exp\Bigl(-\frac{(B_*-\mu)^2}{2m_0}\Bigr)
\]
for a stationary Gaussian approximation, is determined by the spectral
measure of \(B_{ij}(t)\) once the level is given.
A two-mode parent is the comb, not a constant hazard.
The exact codimension-\(2\) condition has rate zero.

Radiative decay of an in-band resonance is a third clock, with width
controlled by \eqref{eq:dos}, and it does not read \(Q\).
You could separate a parent with an empty store and an eigenvalue inside
the band from a parent that is spectrally bound and carries a stored
\(\xi\).
The first clock is the unitary width.
The second, if it splits, replays \(\xi\) and the daughters are the STORE
preimage.
The marginal of \(P_-\) on a Haar pair is Theorem~\ref{thm:haar},
with no free parameter.
A departure from that law, from the arcsine law, and from a point mass at
\(1/2\) would be a different distribution.
Agreement with one of the three is the corresponding ensemble.

\section{Expansion identities that assume a metric}
\label{sec:expand}

If a three-dimensional metric is already expanding and the sheets are
comoving, area per physical volume scales as \(a^{-1}\),
pairwise intersection length as \(a^{-2}\),
and triple-point density as \(a^{-3}\).
The relations \(a\propto\rho_S^{-1}\) and \(a\propto\rho_I^{-1/3}\)
are those two dilutions.
They do not determine \(\ddot a\).
Given an independent \(\ell(t)\propto t^\alpha\),
a single flat-universe fluid would have
\[
q=\frac{1-\alpha}{\alpha},\qquad
w=\frac{2-3\alpha}{3\alpha}.
\]
No earned equation produces \(\alpha\).
In particular \(w=-1\) is not implied by \(\dot\rho_S<0\).

On the chain, the derived spacing is \(\ell_W\propto\lambda^{-1/2}\).
The exponent \(-1/3\) belongs to triple points in the three-volume,
not to the stitch intensity \(\lambda\).
While \(\lambda\) is fixed, Section~\ref{sec:codim} gives \(\dot a=0\).

For a pure power \(V=A(r/\ell)^d\),
\(d_{\mathrm{eff}}=d\) is independent of \(\ell\).
On the chain, \(d_{\mathrm{eff}}\) is a function of \(r/\ell_W\) only:
\(2\), then \(3\), then \(2\) again past \(S\ell_W\).
Walking in \(r\) on a fixed graph traces that curve.
A time-dependent exponent at fixed probe radius needs \(\ell_W(t)\) or \(S(t)\).
The local spectral gap of a tense motif stays fixed in hopping units if
local edges are held fixed while \(\ell_W\) grows.
The probe remains inside the three-dimensional window only while
\(\ell_W\ll r_{\mathrm{probe}}\ll S\ell_W\).

Static response of a drive detuned by \(1/2\) below the \(J_2\) edge is the
massive symbol \(1/2+2|q|^2\), Helmholtz constant \(\kappa=1/2\).
The gapless Laplace response sits at the band edge, which the
electromagnetic falsification labels as excluded tuning.
On a three-dimensional quotient that excluded Green function would be
\(1/r\).
\(J_2\) is \(8\)-regular against the Maxwell value \(4\), so excess
coordination is already present.
Curvature of \eqref{eq:ds} is the curvature of an anisotropic background
whose axis ratio is \(\ell_W\), not a linearization about \(\delta_{ab}\),
and only where \(\ell_W\) varies in space.
A uniform \(\ell_W\) has none of that curvature.
Between events a frozen \(Q\) does not source it, because
\(\partial H/\partial Q=0\).

\section{Plaquettes as an order parameter}
\label{sec:soup}

A tree has exponential ball growth and no integer plateau of
\(d_{\mathrm{eff}}\).
An expander has diameter \(\sim\log N\).
The earned geometric graphs sit between them, with a density of \(4\)-cycles
of order one per vertex.
You could take the fraction \(f_\square\) of edges that lie in a \(4\)-cycle
as a coordinate on that interval:
\(0\) on a tree, order one on the square and on \(J_2\).
Across a percolation of those cycles at fixed degree, a logarithmic
derivative that was large on the tree side falls toward \(2\).
A later weave window at \(3\) is Section~\ref{sec:ell}, and it is
intermediate when \(S\) is finite.
A rise of \(d_{\mathrm{eff}}\) from \(1\) toward \(3\) is a different path
through the same \((f_\square,d_{\mathrm{eff}})\) plane.
